\documentclass{ufctex}

\ufcsetup{
  tipo = projeto,
  impressao = anverso,
  capa = sim,
  ies = {Instituição de Origem Teste},
  entidade-submissao = {Universidade Federal do Ceará},
  programa-projeto = {Programa de Pós-Graduação de Teste},
  tipo-projeto = {Projeto de pesquisa},
  autor = {Autor Projeto Teste},
  titulo = {Projeto de Pesquisa Normativo},
  subtitulo = {Validação da NBR 15287},
  variacao-titulo = {Project Research Normative Test},
  volume = {2},
  local = {Fortaleza},
  ano = {2026},
  orientador = {Prof. Orientador Projeto Teste}
}

\ufcbibliografia{tests/fixtures/referencias-v2.bib}

\begin{document}
\pretextual

\imprimircapa
\imprimirfolhaderosto
\imprimirsumario

\textual
\section{Introdução}
A introdução apresenta o tema, o problema, os objetivos e a justificativa do projeto.

\subsection{Problema de pesquisa}
Problema de pesquisa da fixture normativa.

\subsection{Objetivos}
Objetivos da fixture normativa.

\subsection{Justificativa}
Justificativa da fixture normativa.

\section{Referencial teórico}
Referencial teórico utilizado para embasar a proposta.

\section{Metodologia}
Métodos, procedimentos de coleta e estratégia de análise.

\section{Recursos}
Recursos necessários à execução da pesquisa.

\section{Cronograma}
Cronograma das atividades previstas para a pesquisa.

\nocite{silva2020}
\imprimirreferencias

\end{document}
